\begin{ledgroupsized}[r]{120mm}
\footnotesize
\pstart
\noindent\textbf{Überlieferung:}
\pend
\end{ledgroupsized}
\begin{ledgroupsized}[r]{114mm}
\footnotesize
\pstart \parindent -6mm
\makebox[6mm][l]{\textit{L}}%
Konzept:
LH~XXXVII~5~Bl.~62. 
Ein Blatt~4\textsuperscript{o};
\href{\Scan{LH_37_05_062wz}}%
{Wasserzeichen};
Papiererhaltungsmaßnahmen;
oberer, linker und rechter Rand beschnitten.
Eine Seite auf Bl.~62~r\textsuperscript{o}; Bl.~62~v\textsuperscript{o} leer bis auf die eigh.\ Anmerkung
(in Bleistift, quer): 
\textit{Geometriae Organicae onus} \lbrack/\rbrack\
\textit{id, quo} \textlangle\textit{ser}\textrangle \textit{vit der Geometriae} \lbrack/\rbrack\ 
\textit{pro meis \greek{f}oronomicis} \lbrack/\rbrack\ 
\textit{illustrandis}.
Zum Begriff \glqq geometria organica\grqq\ (d.\,h. geometrische Beweisführung unter Verwendung komplexerer Hilfsmittel als Zirkel und Lineal) siehe bspw. 
\cite{02155}\textit{De geometria seu potius algebra mechanica} von 1673
(\textit{LSB} VII, 1 N.\;8, bes. S.\;104\textendash106).
\pend
\end{ledgroupsized}
\selectlanguage{latin}
\vspace{8mm}
\pstart
\centering
\href{\Scan{LH_37_05_062r}}%
{\lbrack62~r\textsuperscript{o}\rbrack}
\hfill
Accessio calculi 
\edlabel{37_05_062_1a}%
\edtext{}{%
{\xxref%
{37_05_062_1a}{37_05_062_1b}}%
\lemma{dynamici}%
\Bfootnote{%
\textit{(1)}~$av+$
\textit{(2)}~Primo $av+b$
\textit{(3)}~(1)\textsuperscript{mo}
\textit{(4)}~\textit{a}~\textit{L}%
}}%
dynamici
\hfill
\phantom{\lbrack62~r\textsuperscript{o}\rbrack}
\pend
\vspace{0.5\baselineskip}
\pstart 
\begin{tabular}[b]{ccccc}\textit{a}\edlabel{37_05_062_1b}&ante ictum\indexS{ictus}&\textit{v}&post ictum&\textit{y}\\\textit{b}&ante ictum&\textit{x}&post ictum\indexS{ictus}&\textit{z}\end{tabular} %
\pend
\pstart\noindent
\edtext{ponendo \textit{v}, \textit{x}, \textit{y}, \textit{z} esse velocitates sumtas ad easdem partes ad quas maxima est progressio.\indexS{progressio} Quod si corpus contra procedat significatur quantitas negativa.}{%
\lemma{}%
\Bfootnote{%
ponendo \lbrack...\rbrack\  negativa. %
\textit{erg.~L}}}
\pend
\pstart 
\edtext{(Primo)}{%
\lemma{}%
\Bfootnote{%
(Primo) %
\textit{erg.~L}%
}}
$avv+bxx=ayy+bzz$.
(2\textsuperscript{do}) 
$av+bx=%
\edlabel{37_05_062_2a}%
\edtext{}{%
{\xxref%
{37_05_062_2a}{37_05_062_2b}}%
\lemma{}\Bfootnote{%
$ay+bz$.
\bars{Tollatur \textit{a}, \textit{b} per 1 et 2 prodibit aeq.~3 \textit{erg.\ u.\ gestr.}}
\bars{Tollendo $a:b$ prodibit \textit{erg.}} 
(3\textsuperscript{tio})~\textit{L}}}%
ay+bz$. Tollendo $a:b$ prodibit (3\textsuperscript{tio})\edlabel{37_05_062_2b}
$v-x=-y+z$. \pend
\pstart Primum 
\edtext{Theorema}{\lemma{}\Bfootnote{Theorema \textit{erg.~L}}} 
sic 
\edtext{enuntiatur\lbrack:\rbrack\ potentia ante concursum\indexS{concursus}\indexS{potentia ante concursum} est aequalis}{\lemma{enuntiatur}\Bfootnote{%
\textit{(1)}~vires ante ictum sunt aequ %
\textit{(2)}~potentia ante %
\textit{(a)}~ictum %
\textit{(b)}~concursum est aequalis~\textit{L}}} 
potentiae post concursum.\indexS{concursus}\indexS{potentia post concursum}\pend
\pstart Secundum ita: progressio ante concursum\indexS{concursus}\indexS{progressio ante concursum} est aequalis progressioni
\edlabel{37_05_062_3a}%
\edtext{}{{\xxref{37_05_062_3a}{37_05_062_3b}}%
\lemma{post concursum.}%
\Bfootnote{%
\textit{(1)}~Tertium tandem \bars{sic: \textit{streicht Hrsg.}} %
\textit{(a)}~celeritas respectiva, ante ictum, est aequalis %
\textit{(b)}~celeritas mutatae d %
\textit{(2)}~Tertium~\textit{L}}}%
post concursum.\indexS{concursus}\indexS{progressio post concursum}
\pend
\pstart Tertium\edlabel{37_05_062_3b}
ita: celeritas respectiva seu mutandae distantiae\indexS{celeritas mutandae distantiae} ante concursum,\indexS{concursus}\indexS{celeritas respectiva ante concursum} 
aequatur celeritati respectivae post concursum.\indexS{concursus}\indexS{celeritas respectiva post concursum} 
Posset et sic enuntiari 
\edtext{(4\textsuperscript{to})}{\lemma{}\Bfootnote{(4\textsuperscript{to}) \textit{erg.~L}}} 
$v+y=x+z$ seu summa celeritatum progrediendi\indexS{celeritas progrediendi} ejusdem corporis\indexS{summa celeritatum progrediendi ejusdem corporis} ante et post concursum\indexS{concursus} in utroque corpore eadem est. 
Et ex artic.~2.\ fieri potest 
\edtext{(5\textsuperscript{to})}{\lemma{}\Bfootnote{(5\textsuperscript{to}) \textit{erg.~L}}} 
$v-y,:,-x+z=b:a$ id est differentiae celeritatum progrediendi\indexS{celeritas progrediendi} ejusdem corporis\indexS{differentia celeritatum progrediendi ejusdem corporis} sunt reciproce ut duo corpora. 
\pend
\pstart
Unde ex 4 et 5 fit differentiae quadratorum a celeritatibus\indexS{quadratum celeritatis} ejusdem corporis\indexS{differentia quadratorum celeritatum progrediendi ejusdem corporis} sunt reciproce 
\edtext{ut duo corpora, seu}{\lemma{ut}\Bfootnote{%
\textit{(1)}~corpora %
\textit{(2)}~duo corpora, %
\textit{(a)}~seu 6. %
\textit{(b)}~seu %
\textit{(c)}~seu~\textit{L}}} 
(6\textsuperscript{to}) $vv-yy,:,-xx+zz=b:a$. Sed haec aequatio etiam patet ex aequ.~1. %
\pend
\pstart
\edlabel{37_05_062_7a}%
Haec 
\edtext{jam theoremata}{\lemma{jam}\Bfootnote{%
\textit{(1)}~ad %
\textit{(2)}~theorema caeter %
\textit{(3)}~theoremata~\textit{L}}} 
suo quodam modo ad tria etiam et plura corpora produci possunt. Semper enim quantitas virium\indexS{quantitas virium} et quantitas progressionis\indexS{quantitas progressionis} sunt aequales. 
In duobus tamen etiam celeritates respectivae\indexS{celeritas respectiva} sunt aequales. 
Itaque in duobus corporibus elegantissime dici 
\edtext{potest, in duobus}{\lemma{potest,}\Bfootnote{%
\textit{(1)}~conservari %
\textit{(2)}~in duobus~\textit{L}}} 
corporibus concurrentibus conservantur vires,\indexS{vis}\indexS{conservatio virium} progressiones,\indexS{progressio}\indexS{conservatio progressionis} 
\edtext{et respectivae celeritates.\indexS{celeritas respectiva}\indexS{conservatio celeritatis respectivae}}{\lemma{et}\Bfootnote{%
\textit{(1)}~celeritates respective %
\textit{(2)}~respectivae celeritates.~\textit{L}}}
Tria scilicet quorum primum est trium,
secundum duarum,
3tium unius dimensionis.\indexS{dimensio}%
\edlabel{37_05_062_7b}%
\pend
\pstart Jam alias literas assumamus, ad exprimendas
\edtext{ipsas}{%
\lemma{}%
\Bfootnote{%
ipsas %
\textit{erg.~L}%
}}
progressiones. 
\edtext{Sit (7\textsuperscript{mo})}{%
\lemma{Sit}%
\Bfootnote{%
\textit{(1)}~$av=11$ et %
\textit{(2)}~(7\textsuperscript{mo})~\textit{L}%
}}
$av=11$, et $bx=21$ et $ay=12$ et $bz=22$. Ex 2 fiet 
\edlabel{37_05_062_4a}%
\edtext{}{{\xxref{37_05_062_4a}{37_05_062_4b}}%
\lemma{(8\textsuperscript{vo}) \lbrack...\rbrack\  (9\textsuperscript{no})}%
\Cfootnote{Die Zählungen (8\textsuperscript{vo}) und (9\textsuperscript{no}) werden von Leibniz je zweimal verwendet.}}%
(8\textsuperscript{vo}) 
$11+21=%
\edtext{12+22$. Et}{\lemma{$12+22$.}\Bfootnote{%
\textit{(1)}~Et ex 3 fiet %
\textit{(2)}~Et~\textit{L}}} 
per 7.\ fiet (8\textsuperscript{vo}) $v=11:a$ et $x=21:b$ et $y=12:a$, et $z=22:b$. Ergo ex artic.~3 per 8, fiet 
\edtext{(9\textsuperscript{no})}{\lemma{}\Bfootnote{(9\textsuperscript{no}) \textit{erg.~L}}} 
$(11:a)-(21:b)=-(12:a)+(22:b)$ vel $11+12\,,:a=21+22\,,:b$ 
seu (9\textsuperscript{no})\edlabel{37_05_062_4b}
$11+12,:,21+22=a:b$ id est summae progressionum\indexS{progressio}
ejusdem corporis\indexS{summa progressionum ejusdem corporis} sunt ut 
\edlabel{37_05_062_5a}%
\edtext{}{%
{\xxref%
{37_05_062_5a}{37_05_062_5b}}%
\lemma{corpora.}%
\Bfootnote{%
\textit{(1)}~Ex %
\textit{(2)}~Denique ex %
\textit{(3)}~Et ex %
\textit{(4)}~Jam \lbrack...\rbrack\ $-21+22$. %
\textit{(a)}~Ergo fit (11\textsuperscript{mo}) %
\textit{(b)}~Fiet %
\textit{(c)}~Eodem modo %
\textit{(aa)}~fit %
\textit{(bb)}~(10\textsuperscript{mo}) sit $avv=13$ $bxx=23$ $ayy=14$ $bzz=24$. Ergo ex 1 per 10 fit \bars{(11\textsuperscript{mo}) \textit{erg.}} $13+23=14+24$.
\textit{(aaa)}~Et
\textit{(bbb)}~Et (12\textsuperscript{mo}) $v=$  %
\textit{(aaaa)}~$\surd 13:a$ %
\textit{(bbbb)}~$\surd (13:a)$ et %
\textit{(aaaaa)}~(13\textsuperscript{tio}) %
\textit{(bbbbb)}~$x=\surd (23:b)$ et $y=(14:a)$ et $z=\surd (24:b)$ %
\textit{(d)}~Ergo fie %
\textit{(e)}~Ergo ha %
\textit{(f)}~Ergo habebimus (11\textsuperscript{mo})~\textit{L}}}%
corpora. Jam ex 8 fit (10\textsuperscript{mo}) $11-12=-21+22$. %
\pend	
\pstart
Ergo habebimus (11\textsuperscript{mo})\edlabel{37_05_062_5b}
$11^2-12^2,:,-21^2+22^2=a:b$ seu differentiae quadratorum a progressionibus ejusdem corporis,\indexS{differentia quadratorum progressionum ejusdem corporis} sunt ut 
\edtext{corpora. Comparando}{\lemma{corpora.}\Bfootnote{%
\textit{(1)}~Unde ex 11 et 6 sequitur %
\textit{(2)}~Comparando~\textit{L}}}
jam articulos 5 et 9, fiet\lbrack:\rbrack\ summae progressionum 
\edtext{ejusdem\indexS{summa progressionum ejusdem corporis}}{\lemma{}\Bfootnote{ejusdem \textit{erg.~L}}} 
sunt reciproce ut differentiae celeritatum 
\edtext{progrediendi ejusdem\indexS{differentia celeritatum progrediendi ejusdem corporis} computando}{\lemma{}%
\Bfootnote{progrediendi \bars{ejusdem \textit{erg.}} 
\textit{(1)}~tam
\textit{(2)}~computando~\textit{L}}} 
unamquamque ex eo quod in eodem corpore fit ante et post concursum,\indexS{concursus} seu (12\textsuperscript{mo}) 
$v-y,:,-x+z=11+12,:,
\edtext{21+22$. Et conjungendo articulos 6 et 11, fiet: differentiae}{\lemma{$21+22$.}\Bfootnote{%
\textit{(1)}~\bars{Et \textit{streicht Hrsg.}} %
 similiter summae a quadratis %
\textit{(2)}~Et conjungendo articulos 6 et 11, fiet: %
\textit{(a)}~summae a quadrat %
\textit{(b)}~differentiae~\textit{L}}} 
a quadratis progressionum ejusdem\indexS{differentiae quadratorum progressionum ejusdem corporis} sunt 
\edtext{reciproce ut differentiae}{\lemma{reciproce}\Bfootnote{%
\textit{(1)}~ut summae a quadratis %
\textit{(2)}~ut differentiae~\textit{L}}} 
a quadratis celeritatum progressionum ejusdem,\indexS{differentia quadratorum celeritatum progressionum ejusdem corporis} vel quia in quadratis evanescunt differentia signa, 
idemque est quadratum celeritatis\indexS{quadratum celeritatis} et celeritatis progrediendi,\indexS{quadratum celeritatis progrediendi} idemque quadratum quantitatis motus\indexS{quantitas motus}\indexS{quadratum quantitas motus} 
\edtext{seu motionis,}{\lemma{}\Bfootnote{seu motionis \textit{erg.~L}}}
et progressionis;\indexS{quantitas motionis et progressionis} 
\edtext{ideo fient}{\lemma{ideo}\Bfootnote{%
\textit{(1)}~fiet %
\textit{(2)}~fient~\textit{L}}} 
differentiae a quadratis quantitatum motus,\indexS{differentia quadratorum quantitatum motus} reciproce, 
ut differentiae a quadratis celeritatum.\indexS{differentia quadratorum celeritatum}
Calculo autem ex 6 et 11 fiet 
\edtext{(13\textsuperscript{tio})}{\lemma{}\Bfootnote{(13\textsuperscript{tio}) \textit{erg.~L}}} 
\edtext{$v^2-y^2,:,-x^2+z^2=11^2-12^2,:,-21^2+22^2$.}{%
\lemma{$v^2-y^2 \lbrack...\rbrack\ +22^2$}%
\Cfootnote{%
Die rechte Seite der Gleichung heißt eigentlich $-21^2+22^2,:,11^2-12^2$.}}
Denique ex 11 fiet 
\edtext{(14\textsuperscript{to})}{\lemma{}\Bfootnote{(14\textsuperscript{to}) \textit{erg.~L}}} 
$b11^2+a21^2=b12^2+a22^2$.
\pend
\pstart Si eadem methodo pro viribus\indexS{vis} ipsis exprimendis utamur, rursus theoremata 
\edtext{prodibunt. Sit (15\textsuperscript{mo})}{\lemma{prodibunt.}\Bfootnote{%
\textit{(1)}~Sit %
\textit{(2)}~Sit (15\textsuperscript{mo})~\textit{L}}} 
$avv=13$\lbrack,\rbrack\ $bxx=23$, $ayy=14$\lbrack,\rbrack\ $bzz=24$, quibus jungendo 7, fiet 
\edtext{(16\textsuperscript{to})}{\lemma{}\Bfootnote{(16\textsuperscript{to}) \textit{erg.~L}}} 
$11v=13$ et $21x=23$ et $12y=14$ et $22z=24$. 
\edtext{Ergo (17\textsuperscript{mo}) $v=13:11$}{\lemma{Ergo}\Bfootnote{%
\textit{(1)}~$v=13:11$ %
\textit{(2)}~(17\textsuperscript{mo}) $v=13:11$~\textit{L}}} 
et $x=23:21$ et $y=14:12$ et $z=24:22$. Et per 17 ex artic.~3.\ fiet 
\edtext{(18\textsuperscript{vo})}{\lemma{}\Bfootnote{(18\textsuperscript{vo}) \textit{erg.~L}}} 
$(13:11)-(23:21)= 
\edtext{-(14:12)+(24:22)$. Sed et ex 15 sit (19\textsuperscript{no}) $v=\surd,13:a$}{\lemma{$-(14:12)+(24:22)$.\lbrack,\rbrack}\Bfootnote{%
\textit{(1)}~Seu %
\textit{(2)}~Sed nihil hinc magnopere %
\textit{(3)}~Sed potius %
\textit{(4)}~Sed et %
\textit{(a)}~\textit{va} %
\textit{(b)}~ex 15 %
\textit{(aa)}~fit %
\textit{(bb)}~sit \bars{(19\textsuperscript{no}) \textit{erg.}} %
\textit{(aaa)}~$v=$ %
\textit{(bbb)}~$v=\surd,13:a$~\textit{L}}}
$x=\surd,23:b$\lbrack,\rbrack\ $y=\surd,14:a$\lbrack,\rbrack\ $z=\surd,24:b$. 
\edtext{Porro ex}{\lemma{}\Bfootnote{Porro ex \bars{\textit{streicht Hrsg.}} ex~\textit{L}}} 
1 per 15 fit 
\edtext{(20\textsuperscript{mo})}{\lemma{}\Bfootnote{(20\textsuperscript{mo}) \textit{erg.~L}}} 
$13+23=14+24$. %
\pend
\pstart
Alia etiam Theoremata prodibunt, 
\edtext{si pro}{\lemma{si}\Bfootnote{%
\textit{(1)}~ipsas %
\textit{(2)}~pro~\textit{L}}} 
quadratis celeritatum\indexS{quadratum celeritatis} literas sumamus. \pend
\pstart 
\edlabel{37_05_062_6a}%
Notatu dignum quia conservatur quantitas progressionis,\indexS{quantitas progressionis} 
hinc corpus tantum amittere progressionis\indexS{progressio} quantum dat 
\edtext{alteri. Aggressionem\indexS{aggressio}}{%
\lemma{alteri.}%
\Bfootnote{%
\textit{(1)}~Obgressionem %
\textit{(2)}~Aggressionem~\textit{L}}}
\edtext{a.}{%
\lemma{a.}%
\Cfootnote{\textit{autem}}}
possumus vocare contrariam progressionem.\indexS{progressio} 
Paradoxum servari quantitatem progressionis\indexS{quantitas progressionis}\indexS{conservatio quantitatis progressionis} 
omnibus ad quietem redactis. 
\edtext{Nam totalis progressio\indexS{progressio} 
\edtext{\bracks{nulla}}{%
\lemma{}%
\Bfootnote{nonnulla \textit{L ändert Hrsg.}}}
est, cum aequalis progressio\indexS{progressio} aggressioni.\indexS{aggressio}%
\edlabel{37_05_062_6b}}{%
\lemma{Nam \bracks{...} aggressioni}%
\Cfootnote{%
Zur Emendation von "nonnulla" in "nulla" vgl. die gestrichene Variante \textit{(1)} zu \textit{LSB} VIII,~3 N.\;76, S.\;793.18--794.1 (siehe dazu die Einleitung zu N.\;\ref{intro_serv-quant-progressionis}).}}
\pend 
